\begin{textsample}{POS Dim 1 – to – Score 46.00 – to/to\_ef\_28.txt}  \label{ex:f1_pos_271}
GLOSSÁRIO - LÍNGUA PORTUGUESA Anime O anime é tradicionalmente desenhado à mão.
Porém, com o desenvolvimento dos recursos tecnológicos de animação, principalmente a partir da década de 1990, muitos animes passaram a ser produzidos em computadores.
Os temas abordados nos animes são bem variados ( dramas, ficção, terror, aventura, \textbf{psicologia}, romance, comportamento, mitologia, etc. ).
Outra importante característica dos animes \textbf{atuais} é a ocorrência de elementos tecnológicos nos enredos das histórias.
O anime faz muito sucesso no Japão e em vários países do mundo, incluindo o Brasil.
As animações são elaboradas para o cinema, televisão e revistas em quadrinhos.
Disponível em : http://www.suapesquisa.com/o\_que\_e/anime.htm.
Acesso em 05/07/2019.
Audiobooks Como o nome sugere, audiobooks são livros em formato de áudio, também \textbf{conhecidos} como áudiolivros ou livros falados.
Disponível em : https://www.vidasempapel.com.br/audiobooks/.
Acesso em 08/07/2019.
Blogs Blog é \textbf{uma} palavra \textbf{que} resulta da simplificação do termo weblog.
Este, por sua vez, é resultante da justaposição das palavras da língua inglesa web e log.
Web aparece \textbf{aqui} com o significado de rede ( da internet ), \textbf{enquanto} \textbf{que} log é utilizado para designar o registro de atividade ou desempenho regular de algo.
Numa tradução livre, podemos definir blog como \textbf{um} “ diário online ”.
Disponível em : https://www.significados.com.br/blog/.
Acesso em 08/07/2019.
Booktuber É \textbf{uma} expressão para a definição de quem produz algum canal no YouTube \textbf{focado} em livros e literatura.
Essa expressão tem origem americana para designar as pessoas \textbf{que} produzem material para o nicho, \textbf{uma} vez \textbf{que} é a junção de “ book ” ( livro ) e “ tuber ”, sufixo do termo youtuber ( aquele \textbf{que} produz conteúdo na internet ).
Disponível em : https://valerumlivro.com.br/o-que-e-booktuber/.
Acesso em 08/07/2019.
Boxes O box é \textbf{um} apêndice de \textbf{uma} notícia ou reportagem, sempre editado conjuntamente e com função ilustrativa, ou seja, visual, bem como textual, de forma a valorizar o conteúdo jornalístico.
Texto pequeno, inserido em meio de matéria extensa, normalmente composto em fonte diferente, cercado ou colocado entre fios ou com fundo diferenciado, o box não pode ocupar espaço maior do \textbf{que} a matéria jornalística.
Sua construção possibilita a melhor descrição de ambiente ou personagem e permite ao leitor situar diferentes elementos \textbf{que} interferem na informação principal.
O box permite rememorar acontecimentos e resumir opiniões inseridas no texto da matéria.
Pode aparecer através da transcrição de frases, ou conjuntamente com foto de personalidades, especialistas ou pessoas comuns.
Fonte : https://www.trabalhosfeitos.com/ensaios/Jornalismo-e-Estrutura-De-Texto/62852515.html.
Acesso em 14.06.2019.
Ciberpoemas O ciberpoema constrói -se como \textbf{uma} verdadeira “ instalação ”, ainda \textbf{que} compacta e virtual, onde \textbf{uma} multiplicidade de signos é urdida para dar corpo a \textbf{uma} poética híbrida e impura.
Surge assim, a partir desse intenso processo de cruzamento de linguagens, aquilo \textbf{que} alguns estudiosos chamam de blurred genre, ou seja, gênero híbrido, posto \textbf{que}, em certos casos, não é mais possível classificar determinado gênero textual como pertencente a este ou àquele território ( seja estético ou linguístico ), \textbf{uma} vez \textbf{que} o processo de hibridização o coloca, simultaneamente, tanto neste quanto naquele ( outro ) território.
Como se não houvesse mais \textbf{um} único lugar, mas \textbf{um} “ entrelugar ”, \textbf{que} se avizinha de diversos outros espaços semióticos cuja demarcação não existe de modo delimitado e definido, na medida em \textbf{que} são campos de linguagem imprecisos, vagos e nebulosos nos quais semioses de diversas naturezas entrecruzam -se para dar corpo à literatura eletrônica.
O ciberpoema, como \textbf{um} dos herdeiros das vanguardas europeias do início do \textbf{século} XX e do próprio Concretismo, passa a \textbf{exigir} de quem o lê, além de letramentos digitais, \textbf{uma} \textbf{percepção} especial, como aquela \textbf{exigida} para a pintura e a música, posto \textbf{que} a “ leitura ” do ciberpoema \textbf{implica}, sobretudo, \textbf{uma} contemplação do texto, \textbf{uma} audição das estrofes ou até mesmo a utilização de ambas as habilidades ao mesmo tempo, pois é \textbf{oportuno} \textbf{lembrar} \textbf{que} estamos diante de \textbf{um} gênero híbrido e impuro.
A leitura de \textbf{um} ciberpoema requer do leitor a consciência de \textbf{que} alguns daqueles elementos estruturais \textbf{que} constituíam a poesia cânone não mais estarão presentes nessa nova modalidade de gênero \textbf{literário} surgido com o advento das novas tecnologias de informação e comunicação.
Fonte : Revista Fronteira Z – nº 19 – dezembro de 2017.
Gênero digital \textbf{que} utiliza a linguagem eletrônica como forma de comunicação poética, se dá na convergência de texto, som e imagem, encontrando seu ápice na interatividade.
Disponível em : http://comtextodigitalpet.blogspot.com/2012/07/cyberpoesia-literatura-na-era-digital.html.
Acesso em 08/07/2019.
Curadoria de informação \textbf{Um} dos papéis de \textbf{um} curador da informação digital é gerenciar o conteúdo de ambientes digitais.
A Gestão de Conteúdo, \textbf{que} nada mais é \textbf{que} o gerenciamento de informações geradas a partir de dados, foca em captação, ajuste, distribuição e gerenciamento dos conteúdos para apoio ao processo de negócios por meio de ambientes digitais, como portais, intranets etc.
Esses conteúdos podem ser estruturados ou não, são procedentes de sistemas internos e externos, como bancos de imagem, Gerenciamento de Documentos ( GD ), bancos de dados, arquivos nos diretórios das máquinas dos colaboradores e de qualquer outro arquivo digital como som ou vídeo.
É nessa fase \textbf{que} ocorre a curadoria de conteúdo, a edição de conteúdo objetivando a publicação de \textbf{uma} informação nova.
Disponível em : https://periodicos.sbu.unicamp.br/blog/index.php/2018/07/13/curadoria/.
Acesso em 08/07/2019.
Detonado Detonado ( em inglês walkthrough ) é \textbf{um} termo indicativo para resolução de \textbf{tarefas} passo a passo, principalmente de jogos eletrônicos.
Detonados podem ser compostos principalmente por vídeo com narração em áudio, vídeo com legendas de texto, às vezes por apenas texto e imagens estáticas, indicando a solução para as etapas ou fases do jogo.
Disponível em : http://walkthroughuniverse.blogspot.com/p/o-que-e-walkthrough.html.
Acesso em 08/07/2019.
Elementos notacionais Eles \textbf{dizem} respeito ao emprego de determinadas palavras e expressões \textbf{que} frequentemente geram dúvidas entre os usuários do idioma.
Questões notacionais estão relacionadas com a ortografia de alguns termos, sobretudo aqueles \textbf{que} apresentam similaridades fonéticas e ortográficas com outras palavras.
Na modalidade oral, esse tipo de dúvida linguística não faz muita diferença na produção de sentidos do discurso, mas na modalidade escrita a confusão pode acarretar, além do desvio ortográfico, problemas relacionados com a semântica.
Conhecer bem a escrita das palavras \textbf{influencia} positivamente na elaboração do discurso.
Aparte e à parte : Aparte, conjugação do verbo apartar, \textbf{que} significa separar : Não aparte as ovelhas no pasto!
À parte, locução adverbial \textbf{que} significa colocar de \textbf{lado} : Os livros para doação foram colocados à parte.
Disponível em : https://brasilescola.uol.com.br/gramatica/problemas-notacionais-lingua-portuguesa.htm.
Acesso em 08/07/2019.
Cinésica \textbf{Disciplina} \textbf{que} estuda o significado \textbf{expressivo} dos gestos e dos movimentos corporais \textbf{que} acompanham os atos linguísticos ( posturas, expressões faciais, etc. ) ; estudo da linguagem corporal.
Disponível em : https://www.infopedia.pt/dicionarios/lingua-portuguesa/cinésica.
Acesso em 08/07/2019.
Esquete \textbf{Uma} peça de curta duração, geralmente de caráter cômico, produzida para teatro, cinema, rádio ou televisão.
O termo em inglês com o mesmo significado é “ sketch ”.
Disponível em : https://www.significados.com.br/esquete/.
Acesso em 08/07/2019.
Estilização Ação de atribuir estilo, forma própria e diferente, a alguma coisa : estilização do restaurante ; estilização da própria roupa.
Disponível em : https://www.dicio.com.br/estilizacao/.
Acesso em 08/07/2019.
E-zines E-zine ( contração de electronic e fanzine ) ou webzine é \textbf{um} “ fanzine eletrônico ”.
Trata -se de \textbf{uma} publicação periódica, distribuída por e-mail ou \textbf{postada} num site, e \textbf{que} \textbf{foca} \textbf{uma} área específica ( como informática, literatura, música experimental etc. ).
Possui as características de \textbf{uma} revista ( magazine ), mas em vez de usar o formato tradicional de divulgação ( papel ), lança mão do formato eletrônico, seja como \textbf{um} documento \textbf{que} pode ser aberto por \textbf{uma} aplicação específica ( por exemplo, \textbf{um} arquivo de texto, e-book ( PDF ) ou HTML, geralmente com ligações \textbf{que} permitam percorrê -lo em modo de hipertexto ), seja como \textbf{um} executável para \textbf{uma} plataforma específica.
Disponível em : https://pt.wikipedia.org/wiki/E-zine.
Acesso em 08/07/2019.
Fake News Fake News são notícias falsas divulgadas principalmente nas redes sociais.
Os boatos têm informações irreais \textbf{que} apelam para o emocional do leitor/espectador.
Disponível em : https://brasilescola.uol.com.br/curiosidades/o-que-sao-fake-news.htm.
Acesso em 08/07/2019.
Fanclipes Fanclipes são releituras, em sua grande maioria em tom paródico, de videoclipes famosos feitas por fãs.
A ideia é \textbf{que} o fã “ vire ” o artista e “ interprete ” o clipe tal qual o seu ídolo.
Os fanclipes, muitas vezes, se apresentam toscos e visualmente descuidados, mas há alguns \textbf{que} utilizam eficientes recursos de edição para simular a “ situação ” do vídeo original.
Disponível em : http://www.intercom.org.br/papers/nacionais/2009/resumos/R4-0863-1.pdf.
Acesso em 08/07/2019.
Fandom É a junção das palavras ( fan + kingdom ) e significa \textbf{um} grupo de pessoas \textbf{que} são fãs de \textbf{um} certo artista, de \textbf{um} certo jogador, \textbf{um} filme, \textbf{um} cantor, etc.
Essas pessoas usam as redes sociais para se comunicar e espalhar notícia de seus ídolos.
Disponível em : https://qualeagiria.com.br/giria/fandom/.
Acesso em 08/07/2019.
Fanfics Fanfiction ou fanfic é \textbf{uma} narrativa ficcional escrita e divulgada por fãs em blogs, sites e em outras plataformas, \textbf{que} parte da apropriação de produtos midiáticos como filmes, séries, quadrinhos, videogames etc., sem \textbf{que} haja a intenção de ferir direitos autorais ou obter lucros.
Portanto, tem como finalidade a construção de \textbf{um} universo paralelo ao original e também a ampliação do contato dos fãs com as obras \textbf{que} apreciam para limites mais extensos.
Disponível em :.
Acesso em 28/4/2018.
Revista Na Ponta do Lápis, ano XIV, número 31, páginas 14 e 15, julho de 2018.
Fanvídeo Fanvídeos : são vídeos feitos por fãs, \textbf{que} juntam vídeos de vários fãs ou fã clubes de ídolos ; neles constam o nome da cidade e do estado dos fãs.
Os fãs aparecem com camisetas dos ídolos e fazem coisas diferenciadas para eles.
Disponível em : https://ficsandarts.wordpress.com/o-que-e-fanvideo/ Fanzine A palavra “ fanzine ” nasceu da redução fônica da expressão fanatic magazine.
Ela provém da combinação do final do vocábulo “ magazine ”, \textbf{que} tem o sentido de “ revista ”, com o início de “ fanatic ”.
Trata -se de \textbf{um} veículo editado por \textbf{um} fã, seja de graphic novels, obras de ficção científica ou de poemas, músicas, filmes, videogames, entre outras temáticas incorporadas por estas publicações.
Enfim, são elaboradas por admiradores de certo assunto para pessoas \textbf{que} compartilham a mesma paixão.
Eles podem ser peritos neste campo ou \textbf{simples} entusiastas.
As publicações mais profissionais são conhecidas como “ prozines ”.
Em \textbf{um} ou em outro, os temas podem ser enfocados sob diversas formas : contos, poemas, documentários, quadrinhos, entre outros.
Disponível em : https://clubedolivrodesatolep.wordpress.com.
Acesso em 08/07/2019.
Fanpages Fanpage ou Página de fãs é \textbf{uma} página específica dentro do Facebook direcionada para empresas, \textbf{marcas} ou produtos, associações, sindicatos, autônomos, ou seja, qualquer organização com ou sem fins lucrativos \textbf{que} desejem interagir com os seus clientes no Facebook.
Disponível em : https://www.aldabra.com.br/artigo/o-que-e-uma-fanpage.
Acesso em 08/07/2019.
Fotorreportagens, foto-denúncias, reportagens ( multimidiática ) Fotorreportagens : reportagem essencialmente \textbf{baseada} em fotografias, acompanhadas de pequenas legendas.
Disponível em : https://dicionario.priberam.org/fotorreportagem.
Acesso em 08/07/2019.
Foto-denúncias : \textbf{Uma} denúncia, em sentido genérico, é \textbf{uma} tentativa de levar ao conhecimento público ou de alguma autoridade competente \textbf{um} determinado fato ilegal, aguardando alguma possível suscetível punição. \textbf{Uma} denúncia pode ser também \textbf{um} modo para criticar alguém ou alguma coisa, começando \textbf{um} conflito.
A ideia é \textbf{simples} : não é você \textbf{que} denuncia, é a foto ; ela mostra o fato e, como se \textbf{diz}, vale mais \textbf{que} mil palavras.
Disponível em : https://brainly.com.br/tarefa/21518958readmore.
Acesso em 08/07/2019.
Gameplay Gameplay nada mais é \textbf{que} a junção de todas as experiências de \textbf{um} jogador durante a sua interação com os sistemas de \textbf{um} jogo.
Ou melhor...
Cada vez \textbf{que} você joga, você está interagindo com o jogo e com seus personagens.
Ao final de cada jogada você tem \textbf{um} gameplay \textbf{que} representa todas essas interações realizadas durante a sua permanência no jogo.
Para o desenvolvedor é importante saber os conceitos básicos da construção de jogos para criar experiências lúdicas e desafiadoras \textbf{que} resultam em \textbf{uma} “ jogabilidade ” fluida e envolvente.
Disponível em : http://www.formuladejogos.com.br/single-post/2016/06/27/3-Gameplay-O-que-é-Gameplay.
Acesso em 08/07/2019.
Gêneros hipermidiáticos O gênero hipermídia – \textbf{um} gênero textual \textbf{que} surge a partir da “ fusão ” do hipertexto e da multimídia.
Trata -se de \textbf{um} gênero com grande capacidade de difusão de informações, pois, além de trazer o suporte textual no formato digital e de possibilitar a esse suporte o diálogo com outras linguagens verbais e não verbais, ainda traz em seu corpo links diversos \textbf{que} criam \textbf{um} novo tipo de leitura – não mais linear, mas \textbf{uma} leitura dinâmica, interativa, \textbf{que} permite ao leitor ler fragmentos do texto e partir para outros textos \textbf{que} dialogam com a temática anteriormente lida ( SANTAELLA, 2014 ).
É possível \textbf{dizer} \textbf{que} \textbf{uma} \textbf{marca} característica do texto midiático, ou dos gêneros textuais midiáticos, como o hipertexto e a hipermídia, é o fato de estes conterem “ nós ou pontos de intersecção \textbf{que}, ao serem clicados, \textbf{remetem} a conexões não lineares, compondo \textbf{um} percurso de leitura \textbf{que} salta de \textbf{um} ponto a outro de mensagens contidas em documentos distintos, mas interconectados ” ( SANTAELLA, 2014, p. 7 ).
Disponível em : https://repositorio.ufsc.br/bitstream/handle/123456789/94735/299780.
Acesso em 08/07/2019.
Gêneros multissemióticos Gêneros multissemióticos são gêneros compostos por várias linguagens ( modos e semioses : processo de significação e a produção de significados ).
Isso significa \textbf{que} nas salas de aula devemos dar lugar a gêneros \textbf{que} combinam diferentes modalidades, tais como as linguagens verbal ( oral e escrita ), visual, sonora, corporal e digital.
Disponível em : https://www.escrevendoofuturo.org.br/formacao/pergunte-a-olimpia/178/diferencas-e-aplicacoes-entre-multissemiotico-e-multimidiatico.
Acesso em 05/07/2019.
Disponível em : https://ciberduvidas.iscte-iul.pt/consultorio/perguntas/semiose-e-semiotica/34036.
Acesso em 08/07/2019.
Gêneros multimidiáticos Gêneros multimidiáticos colocam em \textbf{cena} a diversidade de mídias, como a TV, o rádio e a internet.
Essa ideia amplia nossa condição de exploração de gêneros discursivos, à medida \textbf{que} podemos dimensionar o ensino de \textbf{uma} variedade de textos \textbf{contemporâneos}, sem abandonar os tradicionalmente contemplados em nosso planejamento, de forma a assegurar \textbf{um} olhar para além do “ impresso/escrito ”.
Disponível em : https://www.escrevendoofuturo.org.br/formacao/pergunte-a-olimpia/178/diferencas-e-aplicacoes-entre-multissemiotico-e-multimidiatico.
Acesso em 08/07/2019.
GIF GIF ( Graphics Interchange Format ou formato de intercâmbio de gráficos ) é \textbf{um} formato de imagem muito usado na Internet, e \textbf{que} foi lançado em 1987 pela CompuServe, para disponibilizar \textbf{um} formato de imagem com cores em substituição do formato RLE, \textbf{que} era apenas preto e branco. \textbf{Uma} das características mais marcantes do GIF é o intercalamento, \textbf{que} armazena as linhas do desenho fora de ordem de modo a possibilitar \textbf{que} \textbf{uma} imagem parcialmente descarregada seja reconhecida antes mesmo de ser totalmente baixada.
Com isso, o usuário pode cancelar o carregamento ao perceber \textbf{que} não é aquilo \textbf{que} queria.
O formato GIF rapidamente se tornou \textbf{popular} por causa da utilização da compressão de dados LZW, \textbf{que}, por ser muito eficiente, permitia \textbf{que} imagens relativamente grandes fossem baixadas num tempo razoável, o \textbf{que} para a época era muito importante, \textbf{uma} vez \textbf{que} a grande maioria das conexões era feita usando modens muito lentos. \textbf{Um} tipo particular de GIF \textbf{bastante} \textbf{conhecido} é o chamado GIF animado.
Ele na verdade é composto de várias imagens do formato GIF, compactadas em \textbf{um} só arquivo.
Essa variante é utilizada para compactar objetos em jogos eletrônicos, para usar como emoticon em mensageiros instantâneos e para enfeitar sites na Internet.
Apesar do formato GIF atualmente ainda ser muito utilizado na web por conta de seu tamanho compacto, ele tem \textbf{uma} paleta limitada de cores — 256 no máximo —, impossibilitando o seu uso prático na compactação de fotografias.
Por causa disso, o formato GIF é utilizado apenas para armazenar ícones e pequenas animações.
Disponível em : https://www.techtudo.com.br/artigos/noticia/2012/04/o-que-e-gif.html.
Acesso em 08/07/2019.
Happening Quase sempre planejada, incorpora -se algum elemento de espontaneidade ou improvisação, \textbf{que} nunca se repete da mesma maneira a cada nova apresentação.
Apesar de ser definida por alguns historiadores como \textbf{um} sinônimo de performance, o happening é diferente porque, além do aspecto de imprevisibilidade, geralmente envolve a participação direta ou indireta do público espectador.
Para o compositor John Cage, os happenings eram “ eventos teatrais espontâneos e sem trama ”.
O termo happening, como \textbf{categoria} artística, foi utilizado pela primeira vez pelo artista Allan Kaprow, em 1959.
Como evento artístico, acontecia em ambientes diversos, geralmente fora de museus e galerias, nunca preparados previamente para esse fim.
Na pop art, artistas como Kaprow e Jim Dine programavam happenings com o intuito de “ tirar a arte das telas e trazê -la para a vida ”.
Robert Rauschenber, em Spring Training ( do inglês, Treino de Primavera ), alugou trinta tartarugas para soltá -las sobre \textbf{um} palco escuro, com lanternas presas nos cascos. \textbf{Enquanto} as tartarugas emitiam luzes em \textbf{direções} aleatórias, o artista perambulava entre elas vestindo calças de jóquei.
No final, sobre pernas-de-pau, Rauschenberg jogou água em \textbf{um} balde de gelo seco preso a sua cintura, levantando nuvens de vapor ao seu redor.
Ao terminar o happening, o artista afirmou : “ As tartarugas foram verdadeiras artistas, não foi? ”.
Disponível em : nsurretosfuriososdesgovernados.blogspot.com/2009/01/o-happening-do-ingles-acontecimento-e.html.
Acesso em 08/07/2019.
Hiperlinks Hiperlink é sinônimo de link.
Hiperlink consiste em links \textbf{que} vão de \textbf{uma} página da Web ou arquivo para outro(a) ; o ponto de partida para os links é denominado de hiperlinks.
História do hiperlink : o termo “ Hiperlink ” foi escrito em 1965 ( ou, eventualmente, 1964 ) por Ted Nelson no início do projeto Xanadu.
Nelson tinha sido \textbf{inspirado} pelo “ Como eu posso pensar ”, \textbf{um} ensaio \textbf{popular} por Vannevar Bush.
No ensaio, Bush descreveu o Memex, onde qualquer \textbf{um} poderia vincular quaisquer duas páginas de informações em \textbf{uma} “ trilha ” de informação relacionada e, em seguida, era movida para frente e para trás entre páginas como se estivessem em \textbf{um} rolo de microfilme único.
A analogia \textbf{contemporânea} mais próxima seria criar \textbf{uma} lista de marcadores para páginas da Web relacionadas e, em seguida, permitir \textbf{que} o usuário virasse para frente e para trás através da lista. \textbf{Uma} série de livros e artigos publicados desde 1964 através de 1980, Nelson transpôs o conceito de Bush de cruzar automatizado no contexto de computador, deixou strings de texto aplicáveis às específicas, em vez de páginas inteiras, generalizou -o de \textbf{uma} máquina local do porte de secretária a \textbf{uma} rede mundial de computadores \textbf{teórica} e defendeu a criação de \textbf{uma} rede de tal.
Entretanto, trabalhando de forma independente, \textbf{uma} equipe liderada por Douglas Engelbart ( com Jeff Rulifson como programador chefe ) foi a primeira a implementar o conceito de hiperlink para rolagem dentro de \textbf{um} único documento ( 1966 ) e \textbf{logo} depois para a ligação entre parágrafos dentro de documentos separados ( 1968 ).
Consulte NLS. \textbf{Um} programa de banco de dados HyperCard foi lançado em 1987 para o Apple Macintosh \textbf{que} permitia hyperlink entre vários tipos de páginas em \textbf{um} documento.
O artigo principal questões jurídicas : aspectos de direitos autorais do hyperlink e elaborando \textbf{enquanto} hyperlink entre páginas da Web é \textbf{uma} característica intrínseca da web, alguns sites objetam estar sendo vinculados a partir de outros Web sites ; alguns afirmaram \textbf{que} a vinculação a eles não é permitida sem permissão.
Disponível em : https://sites.google.com/site/sitesrecord/o-que-e-um-hiperlink Hipertexto O Hipertexto é \textbf{um} conceito associado às tecnologias da informação e \textbf{que} faz referência à escrita eletrônica.
Desde sua origem, o hipertexto vem mudando a noção tradicional de autoria, \textbf{uma} vez \textbf{que} ele contempla diversos textos.
Trata -se, portanto, de \textbf{uma} espécie de obra coletiva, ou seja, apresenta textos dentro de outros, formando assim \textbf{uma} grande rede de informações interativas.
Nesse sentido, sua maior diferença é justamente a forma de escrita e leitura.
Assim, num texto tradicional a leitura segue \textbf{uma} linearidade, \textbf{enquanto} no hipertexto ela é não linear.
Disponível em : https://www.todamateria.com.br/o-que-e-hipertexto/.
Acesso em 08/07/2019.
O hipertexto é, em sua definição, \textbf{uma} forma de escrita e leitura não linear, com blocos de informação ligados a palavras, partes de \textbf{um} texto ou, por exemplo, imagens.
Disponível em : http://educacao.globo.com/portugues/assunto/estudo-do-texto/hipertexto.html.
Acesso em 08/07/2019.
O hipertexto é, em sua definição, \textbf{uma} forma de escrita e leitura não linear, com blocos de informação ligados a palavras, partes de \textbf{um} texto ou, por exemplo, imagens.
Os textos, ao longo da história da \textbf{humanidade}, apresentam -se, em sua maioria, como narrativas retóricas e lineares, ou seja, a narrativa segue \textbf{uma} temporalidade linear, com acontecimentos subsequentes.
Mas nem sempre foi assim, e \textbf{hoje} em dia também não é mais só assim...
Antes de ter a forma \textbf{que} conhecemos atualmente, o livro já foi de tábuas de argila, de rolos de papiro, de folhas de papel costuradas.
Nos últimos tempos, ganhou formato digital e abandonou as páginas de papel.
No mundo \textbf{contemporâneo}, com o excesso de informações, a narrativa ganha \textbf{uma} estrutura hipertextual, com forma de organização em rede, facilitando a interatividade entre textos necessária para a busca da informação com mais rapidez.
Disponível em : http://educacao.globo.com/portugues/assunto/estudo-do-texto/hipertexto.html.
Acesso em 08/07/2019.
Iconografia Define o estudo dos assuntos representados por imagens artísticas, obras de arte, relacionando com as suas fontes e significados.
Disponível em : https://www.significados.com.br/iconografia/.
Acesso em 08/07/2019.
Indoor A mídia indoor é \textbf{uma} tendência do mercado \textbf{que} consiste na sinalização digital realizada em lugares com grande fluxo de pessoas ou locais com espera forçada, como shoppings centers, aeroportos, saguão de hotéis, elevadores, filas de banco, recepções, entre outros.
Esse tipo de publicidade é vendido naturalmente ao público, \textbf{uma} vez \textbf{que} há \textbf{uma} certa receptividade da audiência nas situações de espera, sendo considerada mais efetiva quando comparada às mídias externas, devido à atenção dispersada das pessoas.
Disponível em : https://blog.teclogica.com.br/midia-indoor-o-que-e-e-por-que-investir/.
Acesso em 08/07/2019.
Infográficos Infográfico é \textbf{uma} ferramenta \textbf{que} serve para transmitir informações através do uso de imagens, desenhos e demais elementos visuais gráficos.
Normalmente, o infográfico acompanha \textbf{um} texto, funcionando como \textbf{um} resumo didático e \textbf{simples} do conteúdo escrito.
Por unir texto e imagens, o infográfico atua em duas zonas distintas do cérebro humano : o \textbf{lado} direito, responsável por entender e interpretar figuras ; e o \textbf{lado} esquerdo, \textbf{que} é \textbf{focado} na escrita e no \textbf{raciocínio} lógico.
Assim, os infográficos acabam por simplificar a interpretação dos conteúdos, pois as duas áreas do cérebro atuam em conjunto.
Disponível em : https://www.significados.com.br/infografico/.
Acesso em 08/07/2019.
Inteligibilidade Capacidade de perceber e compreender bem as coisas, dada toda a complexidade e multiplicidade do nosso mundo.
Disponível em : https://www.dicionarioinformal.com.br/inteligibilidade/.
Acesso em 08/07/2019.
Jingle É \textbf{um} termo inglês cujo significado refere -se à música composta para promover \textbf{uma} \textbf{marca} ou \textbf{um} produto em publicidades de rádio ou televisão.
Disponível em : https://www.significados.com.br/jingle/.
Acesso em 08/07/2019.
Link Link é \textbf{uma} palavra em inglês \textbf{que} significa elo, vínculo ou ligação.
No âmbito da informática, a palavra link pode significar hiperligação, ou seja, \textbf{uma} palavra, texto ou imagem \textbf{que}, quando é clicada pelo usuário, o encaminha para outra página na internet, \textbf{que} pode conter outros textos ou imagens.
Em inglês, a palavra link também é usada para determinar \textbf{que} existe \textbf{uma} ligação entre dois elementos.
Disponível em : https://www.significados.com.br/link/.
Acesso em 08/07/2019.
Machinimas Machinimas são vídeos feitos com base em \textbf{cenas} de videogames de absolutamente qualquer gênero e \textbf{que} contam alguma história.
Disponível em : https://www.tecmundo.com.br/curiosidade/4204-machinima-quando-o-cenario-das-historias-e-o-mundo-dos-games.htm.
Acesso em 08/07/2019.
Mancha gráfica É o espaço delimitado de impressão dentro de \textbf{uma} página, onde cai tinta sobre o papel ; fora destes limites, nada pode ser impresso e nenhum elemento pode ultrapassar.
Nos casos em \textbf{que} a mancha ultrapassa as bordas do papel, \textbf{diz} -se \textbf{que} a impressão é sangrada.
A mancha gráfica, elemento \textbf{que} compõe a identidade visual de \textbf{uma} publicação, representa \textbf{um} molde a ser utilizado em todas as páginas para diagramação de textos, imagens e acessórios informativos, \textbf{funcionais} e decorativos e, se dividida em colunas e espaços horizontais, confere flexibilidade ao trabalho de diagramação.
Disponível em : http://www.comuniqueiro.com/dicionario/mancha-grafica.
Acesso em 08/07/2019.
Mashup Mashup significa misturar.
Essas misturas normalmente acontecem com a junção da parte vocal de \textbf{uma} música com a instrumental de outra, criando composições muitas vezes ótimas e outras nem tanto.
Disponível em : https://medium.com/tend\%C3\%AAncias-digitais/mashup-de-m\%C3\%BAsicas-e-efeitos-sonoros-143a6428618d.
Acesso em 08/07/2019.
Meme É \textbf{um} termo grego \textbf{que} significa imitação.
O termo é \textbf{bastante} \textbf{conhecido} e utilizado no “ mundo da internet ”, referindo -se ao fenômeno de “ viralização ” de \textbf{uma} informação, ou seja, qualquer vídeo, imagem, frase, ideia, música e etc., \textbf{que} se espalhe entre vários usuários rapidamente, alcançando muita popularidade.
Disponível em : https://www.significados.com.br/meme/.
Acesso em 08/07/2019.
Mídia Consistem no conjunto dos diversos meios de comunicação, com a finalidade de transmitir informações e conteúdos variados.
O universo midiático abrange \textbf{uma} série de diferentes plataformas \textbf{que} agem como meios para disseminar as informações, como os jornais, revistas, a televisão, o rádio e a internet, por exemplo.
Disponível em : https://www.significados.com.br/midia/.
Acesso em 08/07/2019.
Modalização É \textbf{um} fenômeno discursivo em \textbf{que} \textbf{um} \textbf{sujeito} falante se coloca como fonte de referências pessoais, temporais, espaciais e, ao mesmo tempo, toma \textbf{uma} atitude em relação ao \textbf{que} \textbf{diz} ou ao seu co-enunciador.
Ela pode ser evidenciada nas manifestações escritas e orais da linguagem, nos mais variados contextos.
Exemplos de modalização : Vai chover amanhã.
Disponível em : https://www.dicionarioinformal.com.br/modaliza\%C3\%A7\%C3\%A3o/ Morfossintaxe Em \textbf{uma} análise, pode -se considerar apenas a palavra ( morfologia ) ou sua função ( sintaxe ).
Quando ambas aparecem juntas, surge a morfossintaxe.
A palavra pode ser classificada isoladamente, analisando apenas sua classe gramatical, ou pode ser estudada a partir da função \textbf{que} estabelece dentro da oração.
Se o objeto de estudo é a palavra, tem -se a análise morfológica.
Entretanto, se a busca é por sua função na oração, surge a análise sintática.
Quando a análise ocorre no âmbito da palavra e da frase, ou seja, quando o estudo envolve classe gramatical e função sintática, temos a Morfossintaxe.
Disponível em : https://brasilescola.uol.com.br/o-que-e/portugues/o-que-e-morfossintaxe.htm.
Acesso em 08/07/2019.
Multiletramento O conceito de “ multiletramentos ” se refere à “ multiplicidade cultural das populações ” e “ à multiplicidade semiótica de constituição dos textos ”, por meio dos quais os \textbf{sujeitos} se informam e se comunicam.
Ou seja, o conceito abarca as noções culturais e de multiplicidade de linguagens, em textos impressos, audiovisuais, digitais ou não.
ROJO, Roxane ; MOURA, Eduardo.
Multiletramentos na escola.
São Paulo : Parábola, 2012.
Multimídia O termo “ multimídia ” nasce da junção de duas palavras : “ multi ”, \textbf{que} significa vários, diversos, e “ mídia ”, \textbf{que} vem do latim “ media ”, e significa meios, formas, maneiras.
Em informática, significa a técnica para apresentação de informações \textbf{que} utiliza, simultaneamente, diversos meios de comunicação, mesclando texto, som, imagens fixas e animadas.
Sem os recursos de multimídia no computador, não poderíamos apreciar os cartões virtuais animados, as enciclopédias multimídia, as notícias veiculadas a partir de vídeos, os programas de rádio, os jogos e \textbf{uma} infinidade de atrações \textbf{que} o mundo da informática e internet nos oferece.
Disponível em : http://www.acessasp.sp.gov.br/cadernos/caderno\_09\_01.php.
Acesso em 08/07/2019.
Paralinguística Ramo da linguística \textbf{que} tem a paralinguagem como objeto de estudo.
Disponível em : http://michaelis.uol.com.br/busca? r=0\&f=0\&t=0\&palavra=paralingu\%C3\%ADstica.
Acesso em 08/07/2019.
Paralinguística é o ramo da linguística \textbf{que} estuda línguas \textbf{que} utilizam referências não verbais à comunicação.
O nome deriva da palavra paralanguage, na paralagem do italiano, usado para descrever os elementos em torno do idioma, tais como : altura, volume e tom de fala.
Disponível em : https://educalingo.com/pt/dic-it/paralinguistica.
Acesso em 08/07/2019.
Pastiche O pastiche é \textbf{uma} obra \textbf{literária} ou artística em \textbf{que} se imita o estilo de outros \textbf{autores}, sejam escritores, pintores, músicos ou outros.
É \textbf{uma} técnica \textbf{que} mescla estilos e \textbf{que} pode ser utilizada com sarcasmo, mas não é sua função, como ocorre no caso da paródia, \textbf{que} satiriza ou critica a obra de origem.
Atualmente, o pastiche pode ser \textbf{visto} como \textbf{uma} espécie de colagem ou montagem, tornando -se recortes de vários textos.
Disponível em : https://linkconcursos.com.br/significado-de-pastiche/.
Acesso em 08/07/2019.
Playlist Significa lista de produção.
Para exemplos de playlists musicais e sugestões de atividades para trabalho, \textbf{veja} : “ Produzir \textbf{uma} playlist comentada com a turma ”.
Revista Na Ponta do Lápis, ano XIV, número 31, páginas 14 e 15, julho de 2018.
Podcasts É \textbf{uma} mídia de transmissão de informações, porém a origem da mídia podcast é muito recente e ainda está em seu processo de crescimento, principalmente no Brasil, onde atinge poucas pessoas.
O podcast é \textbf{um} conteúdo de mídia ( geralmente áudio ) transmitido via RSS.
Você pode usar agregadores como iTunes ou Ziepod para PCs, BeyondPod ou PodStore para Android, Wecast ou o nativo Podcasts para iOS e mais \textbf{uma} infinidade de aplicativos para todas as plataformas.
Disponível em : https://mundopodcast.com.br/artigos/o-que-e-podcast/.
Acesso em 08/07/2019.
Political remix O Remix Político, na maioria dos casos, o produtor ou a equipe de produção geralmente tem \textbf{uma} mensagem clara de \textbf{que} deseja se comunicar usando a peça de remix como veículo para entregar essa mensagem a \textbf{um} público específico.
A maioria dos remixes políticos são críticas das estruturas de poder, destacando injustiças ou chamadas à ação, boicotar empresas ou indivíduos ou participar de movimentos de protesto.
Remixes políticos podem ser pôsteres, revistas, spots de rádio/TV, minidocumentários, qualquer tipo de mídia audiovisual.
Disponível em : http://www.criticalremix.com/home/2011/12/13/political-remix/.
Acesso em 08/07/2019.
Redesign É o ato de reformular, \textbf{renovar}, o design de algo.
Quando se propõe o redesign para algo, é comum manter \textbf{uma} conexão com o \textbf{que} já existia previamente, mantendo assim a identificação por parte dos interlocutores.
Todavia, há casos de \textbf{marcas} cujo redesign trouxe \textbf{um} visual novo, o qual pode ter surgido por readequação ao momento, novo \textbf{posicionamento}, novas tendências.
Disponível em : http://oqueeredesign.blogspot.com/p/o-que-e-redesign.html.
Acesso em 08/07/2019.
Sampleamento das músicas Sample, em inglês, significa “ amostra ”.
Quando usamos \textbf{um} sample de guitarra, por exemplo, estamos adicionando à composição \textbf{uma} amostra de sons tocados na guitarra.
O sample não se limita apenas a gravações de instrumentos reais.
Diferente do remix, o sample é como o recorte de algum trecho de determinada música, usado para criar \textbf{uma} nova melodia com \textbf{uma} nova roupagem do som.
Já o remix se utiliza da mesma estrutura sonora da música original e cria \textbf{uma} nova versão da mesma composição.
Disponível em : https://planetamusica.net/voce-sabe-o-que-e-sample/.
Acesso em 08/07/2019.
Semiose Processo de significação e de produção de significados.
Linguística.
Processo capaz de produzir e gerar signos, partindo da premissa de \textbf{que} há \textbf{uma} relação recíproca entre significado e significante.
Disponível em : https://www.significados.com.br/? s=semiose.
Acesso em 08/07/2019.
Semiótica Segundo registros históricos, a semiótica teve sua origem na Grécia Antiga, mas apenas se desenvolveu no começo do \textbf{século} XX, com o trabalho de alguns \textbf{pesquisadores}, como o mestre da linguística e filósofo Ferdinand de Saussure ( 1857–1913 ), e Charles Peirce ( 1839–1914 ), considerado o “ papa da Semiótica ”.
Através da semiótica, somos capazes de interpretar as palavras \textbf{que} formam \textbf{um} texto linguístico e atribuir \textbf{um} significado para as respectivas sequências de palavras, por exemplo.
No caso da linguagem não verbal, os sinais também são dotados de significados específicos, como os sinais de trânsito, os movimentos, os sons, os cheiros, etc.
Disponível em : https://www.significados.com.br/semiotica/.
Acesso em 08/07/2019.
Slams Os slams de poesia, ou poetry slam, são batalhas de poesia falada, em forma de competição, campeonato, mas também é meio \textbf{um} programa de auditório, \textbf{uma} diversão, \textbf{uma} roda, \textbf{um} acontecimento, \textbf{um} encontro, sobretudo.
Resumindo, é \textbf{uma} competição de poesia falada \textbf{que} tem três regras básicas : poemas próprios, de no máximo três minutos, sem acompanhamento musical.
Revista Na Ponta do Lápis, ano XIV, número 32, página 8, dezembro de 2018.
Slides mestres O slide mestre é o slide superior no painel de miniatura, no \textbf{lado} esquerdo da janela.
Os layouts mestres relacionados são exibidos \textbf{logo} abaixo do slide mestre.
Disponível em : https://support.office.com/pt-br/article/o-que-\%C3\%A9-um-slide-mestre-b9abb2a0-7aef-4257-a14e-4329c904da54.
Acesso em 08/07/2019.
Software É \textbf{uma} sequência de instruções escritas para serem interpretadas por \textbf{um} computador com o objetivo de executar \textbf{tarefas} específicas.
Também pode ser definido como os programas \textbf{que} comandam o funcionamento de \textbf{um} computador.
Disponível em : https://www.significados.com.br/software/.
Acesso em 08/07/2019.
Spot de campanha Spot é \textbf{um} fonograma ( \textbf{uma} forma de comunicação em sons, ou seja, cartas e anúncios em falas, e não por escrito ) utilizado como peça publicitária em rádio, feita por \textbf{uma} locução \textbf{simples} ou mista ( duas ou mais vozes ), com ou sem efeitos sonoros e música de fundo.
O spot é, geralmente, utilizado na publicidade quando há muita coisa a ser transmitida em \textbf{uma} só mensagem.
Disponível em : https://pt.wikipedia.org/wiki/Spot Spotify O Spotify é gratuito para celular e tablet.
Com o Spotify, você tem acesso a \textbf{um} universo de músicas e podcasts.
Você pode ouvir \textbf{um} artista, \textbf{um} álbum ou criar \textbf{uma} playlist com suas músicas favoritas.
Disponível em : https://play.google.com/store/apps/details? id=com.spotify.music\&hl=pt\_BR.
Acesso em 08/07/2019.
TDIC Tecnologia digital da informação e da comunicação.
Disponível em : https://educacao.estadao.com.br/blogs/blog-dos-colegios-santa-maria/tdic-no-cotidiano-escolar/.
Acesso em 08/07/2019.
Textos multimodais São textos de linguagens e suportes variados, isto é, textos \textbf{que} misturam diferentes linguagens, como a linguagem verbal e desenhos ou gráficos, por exemplo.
Disponível em : http://portuguescereja.editorasaraiva.com.br/textos-multimodais-leitura-e-producao/.
Acesso em 08/07/2019.
Textualidade A textualidade é o encadeamento \textbf{que} se dá ao longo do texto através da ligação sem

% matched lemmas: aqui, atual, autor, basear, bastante, categoria, cena, conhecido, contemporâneo, direção, disciplina, dizer, enquanto, exigir, expressivo, focar, funcional, hoje, humanidade, implicar, influenciar, inspirar, lado, lembrar, literário, logo, marca, oportuno, percepção, pesquisador, popular, posicionamento, postar, psicologia, que, raciocínio, remeter, renovar, simples, sujeito, século, tarefa, teórico, um, ver
\end{textsample}
